\section{Introdução e uso deste modelo}\label{sec:introducao}

Este documento de referência da Universidade Federal do Ceará (UFC) é, ao mesmo tempo, um exemplo compilável e um guia visual do \gls{template} \texttt{abntexto-ufc}. Em vez de simular uma pesquisa acadêmica fictícia, as próximas seções explicam a estrutura de um trabalho acadêmico, as principais regras de apresentação e de \gls{normalizacao} exercitadas pelo template e os recursos disponíveis para produção de TCCs, monografias de especialização, dissertações, teses e projetos de pesquisa.

O arquivo principal \texttt{main.tex} concentra os metadados institucionais em \texttt{\string\ufcsetup}. As pastas \texttt{frontmatter}, \texttt{chapters} e \texttt{backmatter} separam o conteúdo por etapa documental. A classe \texttt{abntexto-ufc} centraliza a apresentação e usa \texttt{abntexto} como base editorial, preservando a separação entre conteúdo e formatação característica do \LaTeX{} \parencite{knuth,lamport1994}.

\subsection{Como interpretar as orientações deste PDF}

Ao longo do documento são usados quatro tipos de indicação explícita.

\guianormativa{normas ABNT ou referência técnica aplicável}{Identifica uma regra de apresentação cuja fonte pertence à base normativa auditada do projeto. A redação deste guia é uma paráfrase técnica; o texto integral da norma deve ser consultado na fonte autorizada.}

\guiainstitucional{Identifica uma orientação, resolução, instrução normativa ou prática institucional da Universidade Federal do Ceará que complementa a base ABNT quando aplicável ao depósito ou à apresentação do trabalho.}

\guiapolitica{Identifica uma decisão técnica ou editorial do \texttt{abntexto-ufc}. Uma política do modelo não deve ser interpretada automaticamente como exigência universal da ABNT ou da UFC.}

\guiaexemplo{Identifica conteúdo incluído apenas para demonstrar visualmente uma possibilidade do template. Exemplos podem ser removidos ou substituídos no trabalho real.}

Essa separação é importante. Por exemplo, o documento de referência declara PDF/A-2b e a suíte de release o valida com veraPDF; entretanto, o perfil PDF/A-2b específico é uma política técnica do projeto, enquanto a necessidade de PDF/A no fluxo institucional de depósito é tratada separadamente pela documentação da UFC.

\subsection{Base normativa adotada}

A apresentação geral dos trabalhos acadêmicos é organizada a partir da ABNT NBR 14724:2024, considerando a versão corrigida em 1º de abril de 2025. Citações seguem a ABNT NBR 10520:2023 e referências seguem a ABNT NBR 6023:2025 \parencite{abnt14724,abnt10520,abnt6023}.

Outras referências atuam em escopos específicos: ABNT NBR 6028:2021 para resumo e palavras-chave; ABNT NBR 6024:2012 para numeração progressiva; ABNT NBR 6027:2012 para sumário; ABNT NBR 15287:2025 para projetos de pesquisa; ABNT NBR 6034:2004 para índice; ABNT NBR 12225:2023 para lombada; e Normas de Apresentação Tabular do IBGE para tabelas numéricas \parencite{abnt6028,abnt6024,abnt6027,abnt15287,abnt6034,abnt12225,ibge1993}.

\guiainstitucional{O Sistema de Bibliotecas da UFC mantém orientações de normalização e informa que os trabalhos acadêmicos devem observar as normas vigentes aplicáveis. Guias institucionais históricos continuam úteis para requisitos próprios da UFC, mas uma referência antiga a uma edição substituída da ABNT não prevalece sobre a edição vigente adotada pelo projeto \parencite{ufc2026normalizacao,ufcguia2022}.}

\subsection{O que este documento demonstra}

O PDF foi configurado no perfil \texttt{undergraduate-capstone}, em impressão apenas no anverso, com fonte Times em modo portátil, listas, glossário, índice, códigos, algoritmos e tabelas \texttt{tabularray}. Essa configuração serve como corpus visual completo; outros perfis alteram automaticamente os elementos que lhes são aplicáveis.

A suíte de validação do repositório cobre, entre outros pontos, tamanho A4, margens, paginação, tipografia, pré-textuais, citações, referências, objetos, tabelas, equações, códigos, algoritmos, projetos, elementos pós-textuais, matriz de perfis, motores de compilação, compatibilidade com ambiente semelhante ao Overleaf e certificação de Times New Roman e Arial literais em Windows.

\guiapolitica{O template reduz o trabalho manual de formatação, mas não substitui a verificação final das exigências específicas do curso, programa, edital, repositório institucional ou procedimento administrativo aplicável ao trabalho.}

\index{normalização}\index{LaTeX}\index{ABNT}\index{UFC}
